% Source: physical PDF pp. 385--393 (printed pp. 362--370). Coverage:
% complete Section 22.2, from its heading through immediately before
% Section 22.3; Eqs. (22.2.1)--(22.2.49), four unnumbered displays, one
% automatic footnote, one centered divider, citations [6] and [6a], and
% no figures or tables. Source anomalies are itemized in wave1-agent3.md.
\subsection{Transformation of the Measure: The Abelian Anomaly}
\label{sec:22.2}

We now turn to a calculation of anomalies of the sort relevant to $\pi^0$
decay. For this purpose we adopt Fujikawa's
interpretation~\hyperref[ch22-ref-6]{[6]} of the anomaly as a symptom of
the impossibility of defining a suitably invariant measure for integrations
over fermionic field variables. Fujikawa's analysis of this problem was
based on the use of path integrals in Euclidean space, and on an expansion
of the fermionic variables of integration in eigenfunctions of the
gauge-invariant Dirac operator, which is Hermitian in four-dimensional
Euclidean space. Here we shall first present a less rigorous derivation,
based on the familiar path integrals in Minkowski space, which will allow
us to obtain the correct answer with a minimum of work. The Euclidean
approach will be taken up briefly at the end of this section, and used to
derive a famous index theorem.

We will start by evaluating the anomaly in the transformation of the
measure under an arbitrary local matrix transformation
$\psi(x)\to U(x)\psi(x)$ of a column $\psi_n(x)$ of massless complex spin
$1/2$ fermion fields that interact non-chirally with a set of gauge fields
$A_a^\mu(x)$ (such as the $u$ and $d$ quark fields interacting with
the electromagnetic vector potential $A^\mu(x)$ in the problem of
calculating the rate for $\pi^0\to2\gamma$). These are fermionic variables,
so the measure is transformed not with the determinant of the
transformation matrix, but with its inverse:
\begin{equation}
 [d\psi][d\bar\psi]\longrightarrow
 \bigl(\operatorname{Det}\mathcal U\,
       \operatorname{Det}\overline{\mathcal U}\bigr)^{-1}
 [d\psi][d\bar\psi],
 \tag{22.2.1}\label{eq:22.2.1}
\end{equation}
where
\begin{equation}
 \mathcal U_{xn,ym}\equiv U(x)_{nm}\delta^4(x-y),
 \tag{22.2.2}\label{eq:22.2.2}
\end{equation}
\begin{equation}
 \overline{\mathcal U}_{xn,ym}
 \equiv\bigl[\gamma_4U(x)^\dagger\gamma_4\bigr]_{nm}\delta^4(x-y),
 \tag{22.2.3}\label{eq:22.2.3}
\end{equation}
and $\gamma_4\equiv i\gamma^0$ is the matrix used to define
$\bar\psi=\psi^\dagger\gamma_4$. The indices $n,m$ run over flavor labels
and Dirac spin indices.

The reader may wonder at this point why we bother to include factors of
$\gamma_4$ in Eq.~\eqref{eq:22.2.3}, since they just amount to a unitary
transformation that should not affect the determinants. The answer is that
in order for calculations to be meaningful we will find it necessary to
regulate the sum over fermion modes in calculating propagators and
determinants, and we shall find that the $\gamma_4$ factors do affect the
regulated determinants. Whether or not we include the $\gamma_4$ factors
thus depends on the method of regularization that is used to make these
hand-waving manipulations meaningful. We include the factors of $\gamma_4$
because we wish to regulate in such a way as to preserve Lorentz
invariance, and in the cases of interest here it is
$\gamma_4U(x)^\dagger\gamma_4$ rather than $U(x)^\dagger$ that transforms
as a scalar.

First, let us consider the case where $U(x)$ is a unitary non-chiral
transformation
\begin{equation}
 U(x)=\exp[i\alpha(x)t],
 \tag{22.2.4}\label{eq:22.2.4}
\end{equation}
with $t$ an ordinary Hermitian matrix (not involving $\gamma_5$, but not
necessarily traceless) and $\alpha(x)$ an arbitrary real function of $x$.
In this case $\mathcal U$ is pseudounitary:
\begin{equation}
 \overline{\mathcal U}\mathcal U=\mathbf{1},
 \tag{22.2.5}\label{eq:22.2.5}
\end{equation}
so the measure is invariant under this sort of transformation. In
particular, the symmetry under the gauge group itself, where $t$ is one of
the non-chiral generators $t_a$, is not spoiled by any anomalies.

Next, consider a local chiral transformation, with
\begin{equation}
 U(x)=\exp[i\gamma_5\alpha(x)t],
 \tag{22.2.6}\label{eq:22.2.6}
\end{equation}
with $t$ again an ordinary Hermitian matrix and $\alpha(x)$ again an
arbitrary real function of $x$. In this case $\mathcal U$ is
pseudo-Hermitian:
\begin{equation}
 \overline{\mathcal U}=\mathcal U.
 \tag{22.2.7}\label{eq:22.2.7}
\end{equation}
The measure is not invariant under the chiral transformation; rather, we
have
\begin{equation}
 [d\psi][d\bar\psi]\longrightarrow
 \bigl(\operatorname{Det}\mathcal U\bigr)^{-2}[d\psi][d\bar\psi].
 \tag{22.2.8}\label{eq:22.2.8}
\end{equation}

Now let us specialize to the case of an \emph{infinitesimal} local chiral
transformation. Taking $\alpha(x)$ to be infinitesimal in
Eq.~\eqref{eq:22.2.6}, we have here
\begin{equation}
 [\mathcal U-\mathbf{1}]_{nx,my}
 =i\alpha(x)[\gamma_5t]_{nm}\delta^4(x-y).
 \tag{22.2.9}\label{eq:22.2.9}
\end{equation}
Using the identity $\operatorname{Det}M=\exp\operatorname{Tr}\ln M$ and
the limiting formula $\ln(1+x)\to x$ for $x\to0$, the measure now has the
transformation property
\begin{equation}
 [d\psi][d\bar\psi]\longrightarrow
 \exp\left\{i\int d^4x\,\alpha(x)\mathcal A(x)\right\}
 [d\psi][d\bar\psi],
 \tag{22.2.10}\label{eq:22.2.10}
\end{equation}
where $\mathcal A$ is the anomaly function:
\begin{equation}
 \mathcal A(x)\equiv-2\operatorname{Tr}\{\gamma_5t\}\delta^4(x-x),
 \tag{22.2.11}\label{eq:22.2.11}
\end{equation}
with `$\operatorname{Tr}$' here denoting a trace over both Dirac and species
indices. The measure $[d\psi][d\bar\psi]$ appears in the path integral
weighted with a factor $\exp\{i\int d^4x\,\mathcal L(x)\}$, so the
factor\linebreak
$\exp\{i\int d^4x\,\alpha(x)\mathcal A(x)\}$ in the transformation rule
\eqref{eq:22.2.10} for the measure has the same effect as if the Lagrangian
density $\mathcal L(x)$ were not invariant under these transformations, but
instead $\mathcal L(x)\to\mathcal L(x)+\alpha(x)\mathcal A(x)$. Hence when
we use an effective Lagrangian with the fermions integrated out, to take
account of the anomaly we must be sure to include a non-invariant term so
that
\begin{equation}
 \mathcal L_{\mathrm{eff}}(x)\longrightarrow
 \mathcal L_{\mathrm{eff}}(x)+\alpha(x)\mathcal A(x).
 \tag{22.2.12}\label{eq:22.2.12}
\end{equation}
It now remains to calculate the anomaly function $\mathcal A(x)$.

At first sight it does not look like we can get any definite result for the
anomaly. The delta-function is infinite, but the trace vanishes. To do
better, we must introduce a regulator to make $\delta^4(x-x)$ meaningful.
We can do this in a gauge-invariant manner by inserting a differential
operator $f(-\sl{D}_x^2/M^2)$ acting on the delta-function before taking
its argument to zero:
\begin{equation}
 \mathcal A(x)=-2\left[
  \operatorname{Tr}\left\{\gamma_5t
  f(-\sl{D}_x^2/M^2)\right\}\delta^4(x-y)
 \right]_{y\to x}.
 \tag{22.2.13}\label{eq:22.2.13}
\end{equation}
Here $D_x$ is the Dirac differential operator in the presence of the gauge
field $A_a^\mu(x)$:
\begin{equation}
 (D_x)_\mu\equiv\frac{\partial}{\partial x^\mu}
 -it_a A_{a\mu}(x).
 \tag{22.2.14}\label{eq:22.2.14}
\end{equation}
Also, $M$ is some large mass, eventually to be taken to infinity, and
$f(s)$ is a smooth function, subject only to the condition\footnote{For
instance, we might take $f(s)=e^{-s^2}$, as done originally by Fujikawa,
or $f(s)=1/(1+s)$.} that as $s$ goes from 0 to $\infty$, $f(s)$ must drop
smoothly from 1 to 0:
\begin{equation}
 f(0)=1,\qquad f(\infty)=0,
 \tag{22.2.15}\label{eq:22.2.15}
\end{equation}
\begin{equation}
 sf'(s)=0\quad\text{at}\quad s=0\ \text{and}\ s=\infty.
 \tag{22.2.16}\label{eq:22.2.16}
\end{equation}
Note that we do not take the regulator function to be a function of
$\sl{\partial}$, because we want to preserve gauge invariance, and we do
not take it to be a function of $D^\mu D_\mu$, because we need it to
regulate not only the determinant but also the fermion propagator
$\sl{D}^{-1}$.

To evaluate the expression \eqref{eq:22.2.13}, we use the Fourier
representation of the delta-function, and write the anomaly function as
\[
 \begin{aligned}
 \mathcal A(x)
 &=-2\int\frac{d^4k}{(2\pi)^4}
 \left[
  \operatorname{Tr}\left\{\gamma_5t f(-\sl{D}_x^2/M^2)\right\}
  e^{ik\cdot(x-y)}
 \right]_{y=x}\\
 &=-2\int\frac{d^4k}{(2\pi)^4}
 \operatorname{Tr}\left\{\gamma_5t
 f\left(-[i\sl{k}+\sl{D}_x]^2/M^2\right)\right\}.
 \end{aligned}
\]
(A derivative with respect to $x$ gives zero when acting on the extreme
right in the second expression, but not when acting on
$A_{a\mu}(x)$.) Rescaling the momentum $k^\mu$ by a factor $M$, this
is
\begin{equation}
 \mathcal A(x)=-2M^4\int\frac{d^4k}{(2\pi)^4}
 \operatorname{Tr}\left\{\gamma_5t
 f\left(-[i\sl{k}+\sl{D}_x/M]^2\right)\right\}.
 \tag{22.2.17}\label{eq:22.2.17}
\end{equation}
The argument of the cut-off function may be written
\begin{equation}
 -\left[i\sl{k}+\frac{\sl{D}_x}{M}\right]^2
 =k^2-\frac{ik\cdot D_x}{M}
 -\left(\frac{\sl{D}_x}{M}\right)^2.
 \tag{22.2.18}\label{eq:22.2.18}
\end{equation}
Eq.~\eqref{eq:22.2.17} receives contributions in the limit
$M\to\infty$ only from terms in the expansion of
$f(-[i\sl{k}+\sl{D}_x/M]^2)$ that have no more than four factors of
$1/M$, and also only from terms that contain at least four Dirac gamma
matrices, because otherwise the trace over Dirac indices vanishes. This
leaves us with only the terms of second order in $\sl{D}_x^2$:
\begin{equation}
 \mathcal A(x)=-\int\frac{d^4k}{(2\pi)^4}f''(k^2)
 \operatorname{Tr}\{\gamma_5t\sl{D}_x^4\},
 \tag{22.2.19}\label{eq:22.2.19}
\end{equation}
which are now independent of the regulator mass $M$.

To evaluate the integral over $k$, we perform the same sort of rotation of
the $k^0$ contour of integration as in evaluating Feynman diagrams, so that
in effect $k^0$ is replaced with $ik^4$, with $k^4$ running from $-\infty$
to $+\infty$. (This is a step that can really only be justified by working
with Euclidean path integrals from the beginning.) The integral is then in
effect
\begin{equation}
 \int d^4k\,f''(k^2)
 =i\int_0^\infty 2\pi^2\kappa^3\,d\kappa\,f''(\kappa^2).
 \tag{22.2.20}\label{eq:22.2.20}
\end{equation}
By repeated integration by parts using Eq.~\eqref{eq:22.2.16} and then
Eq.~\eqref{eq:22.2.15}, this is
\begin{equation}
 \int d^4k\,f''(k^2)
 =i\pi^2\int_0^\infty ds\,s f''(s)
 =-i\pi^2\int_0^\infty ds\,f'(s)
 =i\pi^2.
 \tag{22.2.21}\label{eq:22.2.21}
\end{equation}
To calculate the trace, we write
\begin{equation}
 \begin{aligned}
 \sl{D}_x^2
 &=\frac14\{(D_x)^\mu,(D_x)^\nu\}\{\gamma_\mu,\gamma_\nu\}
 +\frac14[(D_x)^\mu,(D_x)^\nu][\gamma_\mu,\gamma_\nu]\\
 &=D_x^2-\frac14it_a F_a^{\mu\nu}
 [\gamma_\mu,\gamma_\nu].
 \end{aligned}
 \tag{22.2.22}\label{eq:22.2.22}
\end{equation}
The only term in $\sl{D}_x^4$ that contributes to the trace over Dirac
indices is the one involving a product of four Dirac matrices, which gives
\begin{equation}
 \operatorname{tr}_{D}\{\gamma_5
 [\gamma_\mu,\gamma_\nu][\gamma_\rho,\gamma_\sigma]\}
 =16i\epsilon_{\mu\nu\rho\sigma},
 \tag{22.2.23}\label{eq:22.2.23}
\end{equation}
where `$\operatorname{tr}_{D}$' denotes a trace over Dirac indices only,
and as usual $\epsilon_{\mu\nu\rho\sigma}$ is the totally antisymmetric
tensor with $\epsilon^{0123}=+1$. Using
Eqs.~\eqref{eq:22.2.21}--\eqref{eq:22.2.23} in Eq.~\eqref{eq:22.2.19}
then gives the anomaly function as
\begin{equation}
 \mathcal A(x)=-\frac{1}{16\pi^2}\epsilon_{\mu\nu\rho\sigma}
 F_a^{\mu\nu}(x)F_b^{\rho\sigma}(x)
 \operatorname{tr}\{t_a t_b t\},
 \tag{22.2.24}\label{eq:22.2.24}
\end{equation}
with `$\operatorname{tr}$' here denoting a trace only over indices
labelling the various fermion species. In the special case in which $t$ is
the unit matrix, the quantity \eqref{eq:22.2.24} is known as the
\emph{Chern--Pontryagin density}.

This result can be expressed in terms of the current associated with the
anomalous symmetry. If we assume for simplicity that the action itself is
invariant under the symmetry transformation
$\psi(x)\to\psi(x)+it\gamma_5\alpha\psi(x)$ with constant infinitesimal
parameter $\alpha$, then as discussed in
\hyperref[sec:7.3]{Section 7.3}, the change in the action when we make such
a transformation with a spacetime-\emph{dependent} parameter $\alpha(x)$
may be written as
$\delta I=\int d^4x\,J_5^\mu(x)\partial_\mu\alpha(x)$, where $J_5^\mu(x)$
is the current that becomes conserved when the field operators satisfy the
dynamical equations that make the action stationary with respect to
arbitrary variations in the fields. When we make the change of variables
$\delta\psi(x)=it\gamma_5\alpha(x)\psi(x)$, the change in the path integral
over the fermionic fields is
\begin{equation}
 \begin{aligned}
 \delta\int[d\psi][d\bar\psi]e^{iI}
 =i\int d^4x\int[d\psi][d\bar\psi]\,
 \bigl[\mathcal A(x)\alpha(x)
 +J_5^\mu(x)\partial_\mu\alpha(x)\bigr]e^{iI}.
 \end{aligned}
 \tag{22.2.25}\label{eq:22.2.25}
\end{equation}
But this is a mere change of variables, and so for arbitrary $\alpha(x)$
it cannot affect the path integral. Therefore for arbitrary gauge fields
\begin{equation}
 \left\langle\partial_\mu J_5^\mu\right\rangle_A
 =\mathcal A
 =-\frac{1}{16\pi^2}\epsilon_{\mu\nu\rho\sigma}
 F_a^{\mu\nu}F_b^{\rho\sigma}
 \operatorname{tr}\{t_a t_b t\},
 \tag{22.2.26}\label{eq:22.2.26}
\end{equation}
where for any operator $\mathcal O$, $\langle\mathcal O\rangle_A$ is the
quantum average of $\mathcal O$ in a fixed background field $A^\mu(x)$:
\begin{equation}
 \langle\mathcal O\rangle_A\equiv
 \frac{\int[d\psi][d\bar\psi]e^{iI}\mathcal O}
 {\int[d\psi][d\bar\psi]e^{iI}}.
 \tag{22.2.27}\label{eq:22.2.27}
\end{equation}

Incidentally, it is possible to rewrite Eq.~\eqref{eq:22.2.26} as a
conservation condition. Consider the special case where
$\operatorname{tr}\{t_a t_b t\}$ is proportional to
$\delta_{ab}$:
\begin{equation}
 \operatorname{tr}\{t_a t_b t\}=N\delta_{ab}.
 \tag{22.2.28}\label{eq:22.2.28}
\end{equation}
Define a current known as the \emph{Chern--Simons class}:
\begin{equation}
 \begin{aligned}
 G^\mu
 &\equiv2\epsilon^{\mu\nu\lambda\rho}
 \left[
 A_{c\nu}\partial_\lambda A_{c\rho}
 +\frac13C_{abc}
 A_{a\nu}A_{b\lambda}A_{c\rho}
 \right]\\
 &=\epsilon^{\mu\nu\lambda\rho}
 \left[
 A_{c\nu}F_{c\lambda\rho}
 -\frac13C_{abc}
 A_{a\nu}A_{b\lambda}A_{c\rho}
 \right],
 \end{aligned}
 \tag{22.2.29}\label{eq:22.2.29}
\end{equation}
which satisfies the identity
\begin{equation}
 \partial_\mu G^\mu
 =\frac12\epsilon^{\mu\nu\lambda\rho}
 F_{c\mu\nu}F_{c\lambda\rho}.
 \tag{22.2.30}\label{eq:22.2.30}
\end{equation}
Eq.~\eqref{eq:22.2.30} allows us to rewrite Eq.~\eqref{eq:22.2.26} as the
condition
\begin{equation}
 \partial_\mu K^\mu=0,
 \tag{22.2.31}\label{eq:22.2.31}
\end{equation}
where
\begin{equation}
 K^\mu\equiv\left\langle J_5^\mu\right\rangle_A
 +\frac{N}{8\pi^2}G^\mu.
 \tag{22.2.32}\label{eq:22.2.32}
\end{equation}

However, we cannot use the conservation of the current $K^\mu$ to argue,
as we did in the previous section, that the process $\pi^0\to2\gamma$ is
suppressed, because in making that argument we assumed not only the chiral
symmetry associated with axial-vector current conservation but also
electromagnetic gauge invariance, and as we see from
Eq.~\eqref{eq:22.2.29}, the current $K^\mu$ though conserved is not
gauge-invariant.

Our derivation of the formula \eqref{eq:22.2.24} for the anomaly shows that
if we had evaluated the anomaly function using a differential operator
$f(-\sl{\partial}_x^2/M^2)$ in place of $f(-\sl{D}_x^2/M^2)$ in
Eq.~\eqref{eq:22.2.13}, we would have obtained a vanishing anomaly
function. In fact, with this regulator procedure the axial-vector current
is $K^\mu$, not $J_5^\mu$. The trouble with this procedure, as mentioned
above, is that the regulator operator is no longer gauge-invariant,
resulting in the presence of a non-gauge-invariant term in $K^\mu$. There
is no regulator procedure for fermion propagators as well as determinants
that is both gauge- and chiral-invariant.

We may now return to the problem that gave anomalies their start, and use
our results to calculate the actual rate of the process $\pi^0\to2\gamma$.
The symmetry of interest here is generated by charge-neutral chiral
transformations of the light quark fields:
\begin{equation}
 \delta u=i\alpha\gamma_5u,\qquad
 \delta d=-i\alpha\gamma_5d.
 \tag{22.2.33}\label{eq:22.2.33}
\end{equation}
This symmetry is anomaly-free in pure quantum chromodynamics, because the
$u$ and $d$ belong to the same representation of the color gauge group, so
that their contributions to the gluon--gluon terms in the anomaly of the
symmetry \eqref{eq:22.2.33} cancel. On the other hand, in the presence of
an electromagnetic field $A^\mu(x)$, this symmetry has an anomaly
\[
 \mathcal A(x)=-\frac{1}{16\pi^2}\epsilon_{\mu\nu\rho\sigma}
 F^{\mu\nu}(x)F^{\rho\sigma}(x)
 \operatorname{tr}\{q^2\tau_3\},
\]
where $q$ is the quark charge matrix and $\tau_3$ is the diagonal
$2\times2$ matrix with elements $+1$ for $u$ and $-1$ for $d$. If we
assume as usual that there are $N_c$ $u$ quarks of charge $2e/3$ and an
equal number of $d$ quarks of charge $-e/3$, the trace is
\[
 \operatorname{tr}\{q^2\tau_3\}
 =N_c\left(\frac{2e}{3}\right)^2(+1)
 +N_c\left(\frac{-e}{3}\right)^2(-1)
 =\frac{N_ce^2}{3},
\]
so the anomaly function here is
\begin{equation}
 \mathcal A(x)=-\frac{N_ce^2}{48\pi^2}
 \epsilon_{\mu\nu\rho\sigma}F^{\mu\nu}(x)F^{\rho\sigma}(x).
 \tag{22.2.34}\label{eq:22.2.34}
\end{equation}
We must then include in the effective Lagrangian terms that under the
chiral transformation \eqref{eq:22.2.33} give it the transformation
\eqref{eq:22.2.12}, that is,
\begin{equation}
 \delta\mathcal L_{\mathrm{eff}}(x)=\alpha\mathcal A(x)
 =-\frac{N_ce^2}{48\pi^2}\epsilon_{\mu\nu\rho\sigma}
 F^{\mu\nu}(x)F^{\rho\sigma}(x)\alpha.
 \tag{22.2.35}\label{eq:22.2.35}
\end{equation}

Under the transformation \eqref{eq:22.2.33} the pion field transforms as
\begin{equation}
 \delta\pi^0=\alpha F_\pi,
 \tag{22.2.36}\label{eq:22.2.36}
\end{equation}
where $F_\pi=184\ \mathrm{MeV}$ is the pion decay amplitude introduced in
\hyperref[sec:19.4]{Chapter 19}. (The normalization convention for this
constant is fixed by our definition of the symmetry generator as
$\gamma_5\tau_3=2\gamma_5t_3$.) It follows that we must include in the
effective Lagrangian a term
\begin{equation}
 \frac{\pi^0(x)\mathcal A(x)}{F_\pi}
 =-\frac{N_ce^2}{48\pi^2F_\pi}\epsilon_{\mu\nu\rho\sigma}
 F^{\mu\nu}(x)F^{\rho\sigma}(x)\pi^0(x).
 \tag{22.2.37}\label{eq:22.2.37}
\end{equation}
Comparing this with the general formula \eqref{eq:22.1.1} for the
effective Lagrangian for the decay $\pi^0\to2\gamma$, we see that the
constant $g$ in \hyperref[eq:22.1.1]{Eq.~(22.1.1)} must have the value
\begin{equation}
 g=\frac{N_ce^2}{48\pi^2F_\pi}.
 \tag{22.2.38}\label{eq:22.2.38}
\end{equation}
(This shows that our previous crude order-of-magnitude estimate
\eqref{eq:22.1.3} was too large by just a factor $6/N_c$.) The rate
\eqref{eq:22.1.2} for pion decay is thus predicted to be
\begin{equation}
 \Gamma(\pi^0\to2\gamma)
 =\frac{N_c^2\alpha^2m_\pi^3}{144\pi^3F_\pi^2}
 =\left(\frac{N_c}{3}\right)^2
 \cdot1.11\cdot10^{16}\ \mathrm{s}^{-1}.
 \tag{22.2.39}\label{eq:22.2.39}
\end{equation}
The observed rate is
$\Gamma(\pi^0\to2\gamma)=(1.19\mathbin{\pm}0.08)\cdot10^{16}\
\mathrm{s}^{-1}$, in good agreement with the theoretical result
\eqref{eq:22.2.39} if and only if $N_c=3$. The success of this calculation
was one of the first pieces of firm evidence that there are three colors of
quarks.

Remarkably, as we have seen in the previous section, shortly after the
discovery of the $\pi^0$ Steinberger calculated $g$ from a single
proton-loop diagram, and obtained the result
$g=e^2G/32\pi^2m_N$, where $G$ is the pseudoscalar pion--nucleon coupling
constant. This result would precisely agree with Eq.~\eqref{eq:22.2.38}
for $N_c=3$ if we used the Goldberger--Treiman relation with $g_A=1$ to
set $G=2m_N/F_\pi$. The correct result is larger than the Steinberger
result by a factor $g_A^2=1.56$. The reason that Steinberger got nearly the
right result is because the answer is determined by the triangle anomaly,
which is proportional to $\operatorname{tr}\{q^2\tau_3\}$. For one proton,
this trace has the value $e^2$, which is the same as the quark value found
above for three colors of quarks.

\begin{center}
 $\ast\qquad\ast\qquad\ast$
\end{center}

As mentioned earlier, a more rigorous derivation of the anomaly can be
given by using path integrals in Euclidean spacetime. (The use of Euclidean
path integrals is briefly discussed in
\hyperref[app:23.A]{Appendix A of Chapter 23}.) We introduce a Euclidean
fourth coordinate $x_4=ix^0=-ix_0$, and correspondingly
$\partial_4=-i\partial_0$, $\gamma_4=i\gamma^0$, and
$A_{4a}=iA_a^0$. The spacetime volume is then expressed as
$d^4x=-i(d^4x)_E$, where $(d^4x)_E$ is the Euclidean volume element
$(d^4x)_E=dx_1dx_2dx_3dx_4$. In Euclidean spacetime the fields $\psi(x)$
and $\bar\psi(x)$ must be regarded as entirely independent, with local
chiral transformations defined by
$\delta\psi(x)=i\alpha(x)t\gamma_5\psi(x)$,
$\delta\bar\psi(x)=-i\alpha(x)\bar\psi(x)t\gamma_5$. The transformation of
the measure is then again given by Eq.~\eqref{eq:22.2.10}, with the anomaly
function $\mathcal A(x)$ given by Eq.~\eqref{eq:22.2.11}. Introducing a
regulator function as before, this again gives the formula
\hyperref[eq:22.2.13]{(22.2.13)} for $\mathcal A(x)$. One great advantage
of the Euclidean approach is that with $x_4$ and $A_{4a}$ real, the
Dirac operator $i\sl{D}$ in Eq.~\eqref{eq:22.2.13} is Hermitian:
\begin{equation}
 i\sl{D}=[i\partial_i+t_a A_{ia}]\gamma_i,
 \tag{22.2.40}\label{eq:22.2.40}
\end{equation}
with $i,j$, etc. here summed over the values 1, 2, 3, 4. It therefore has
orthonormal spinor eigenfunctions $\varphi_\kappa(x)$:
\begin{equation}
 i\sl{D}\varphi_\kappa=\lambda_\kappa\varphi_\kappa,
 \tag{22.2.41}\label{eq:22.2.41}
\end{equation}
\begin{equation}
 \int(d^4x)_E\,\varphi_\kappa(x)^\dagger\varphi_{\kappa'}(x)
 =\delta_{\kappa\kappa'},
 \tag{22.2.42}\label{eq:22.2.42}
\end{equation}
with real eigenvalues $\lambda_\kappa$. We also are assuming, as throughout
this section, that $t$ commutes with $i\sl{D}$, so that we can choose the
$\varphi_\kappa$ so that
$t\varphi_\kappa=t_\kappa\varphi_\kappa$. These eigenfunctions satisfy a
completeness relation
\begin{equation}
 \sum_\kappa\varphi_\kappa(x)\varphi_\kappa^\dagger(y)
 =\delta^4(x-y)\mathbf{1},
 \tag{22.2.43}\label{eq:22.2.43}
\end{equation}
where `$\mathbf{1}$' is the $4\times4$ unit matrix. Therefore the anomaly
function may now be written as the limit of a manifestly convergent sum:
\begin{equation}
 \begin{aligned}
 \mathcal A(x)
 &=-2\lim_{M\to\infty}\operatorname{Tr}
 \left\{\gamma_5t f(-\sl{D}^2/M^2)
 \sum_\kappa\varphi_\kappa(x)\varphi_\kappa^\dagger(x)\right\}\\
 &=-2\lim_{M\to\infty}\sum_\kappa
 t_\kappa f(\lambda_\kappa^2/M^2)
 \bigl(\varphi_\kappa^\dagger(x)\gamma_5\varphi_\kappa(x)\bigr).
 \end{aligned}
 \tag{22.2.44}\label{eq:22.2.44}
\end{equation}

In just the same way that we derived formula \eqref{eq:22.2.24} for the
anomaly function, we can show here that
\begin{equation}
 \mathcal A(x)=\frac{1}{16\pi^2}\epsilon^E_{ijkl}
 F_{ija}F_{klb}\operatorname{tr}\{t_a t_b t\},
 \tag{22.2.45}\label{eq:22.2.45}
\end{equation}
where $\epsilon^E_{ijkl}$ is the totally antisymmetric tensor with
$\epsilon^E_{1234}=+1$. (The difference of sign in
Eqs.~\eqref{eq:22.2.24} and \eqref{eq:22.2.45} arises because
Eq.~\eqref{eq:22.2.45} is missing two factors of $i$ as compared with
Eq.~\eqref{eq:22.2.24}: one factor $i$ from Eq.~\eqref{eq:22.2.23},
because
\[
 \operatorname{tr}_{D}\{\gamma_5
 [\gamma_i,\gamma_j][\gamma_k,\gamma_l]\}
 =16\epsilon^E_{ijkl},
\]
together with another factor $i$ from the replacement of $d^4k$ with
$(d^4k)_E$ in Eq.~\eqref{eq:22.2.20}.)

Now, given any eigenfunction $\varphi_\kappa(x)$ of $i\sl{D}$ and $t$ with
eigenvalue $\lambda_\kappa\ne0$, there is another normalized eigenfunction
$\varphi_{\kappa^-}(x)$ with eigenvalues
$\lambda_{\kappa^-}=-\lambda_\kappa$ and $t_\kappa$, given by
$\varphi_{\kappa^-}(x)=\gamma_5\varphi_\kappa(x)$. (Recall that in the
notation used throughout this book, $\gamma_5$ is the Hermitian matrix
$\gamma_5=-i\gamma_1\gamma_2\gamma_3\gamma_0
=\gamma_1\gamma_2\gamma_3\gamma_4$.) Furthermore
$\bigl(\varphi_\kappa^\dagger(x)\varphi_\kappa(x)\bigr)
=\bigl(\varphi_{\kappa^-}^\dagger(x)\varphi_{\kappa^-}(x)\bigr)$, so the
terms $\kappa$ and $\kappa^-$ cancel in the sum in
Eq.~\eqref{eq:22.2.44}. This leaves just a sum over the eigenfunctions with
$\lambda_\kappa=0$. These are not generally paired; rather, since
$\gamma_5$ anticommutes with $i\sl{D}$, they can be chosen as simultaneous
orthonormal eigenfunctions $\varphi_u,\varphi_v$ of $i\sl{D}$ with
eigenvalue zero \emph{and} of $\gamma_5$ with eigenvalues $+1$ and $-1$,
respectively.
\begin{equation}
 \begin{aligned}
 i\sl{D}\varphi_u&=0,&\qquad \gamma_5\varphi_u&=\varphi_u,\\
 i\sl{D}\varphi_v&=0,& \gamma_5\varphi_v&=-\varphi_v.
 \end{aligned}
 \tag{22.2.46}\label{eq:22.2.46}
\end{equation}
Using the fact that $f(0)=1$,
\hyperref[eq:22.2.44]{Eq.~(22.2.44)} then becomes
\begin{equation}
 \mathcal A(x)=-2\left[
 \sum_u t_u\bigl(\varphi_u^\dagger(x)\varphi_u(x)\bigr)
 -\sum_v t_v\bigl(\varphi_v^\dagger(x)\varphi_v(x)\bigr)
 \right].
 \tag{22.2.47}\label{eq:22.2.47}
\end{equation}
Now, since the $\varphi_u$ and $\varphi_v$ are normalized as in
Eq.~\eqref{eq:22.2.42}, the integral of Eq.~\eqref{eq:22.2.47} gives
\begin{equation}
 \int(d^4x)_E\,\mathcal A(x)
 =-2\left[\sum_u t_u-\sum_v t_v\right],
 \tag{22.2.48}\label{eq:22.2.48}
\end{equation}
with the sums over $u$ and $v$ running over left- and right-handed zero
modes of the operator $i\sl{D}$, respectively. In particular, for the case
where $t$ is the unit matrix, by using Eq.~\eqref{eq:22.2.45} this may be
expressed as a relation between a functional of the gauge field and the
numbers of the zero modes of the Dirac operator with definite chiralities:
\begin{equation}
 -\frac{1}{32\pi^2}\int(d^4x)_E\,\epsilon^E_{ijkl}
 F_{a ij}F_{b kl}\operatorname{tr}[t_a t_b]
 =n_+-n_-,
 \tag{22.2.49}\label{eq:22.2.49}
\end{equation}
where here $n_\pm$ are the numbers of zero modes of $\sl{D}$ that have
eigenvalues $\pm1$ for $\gamma_5$. This is the celebrated
\emph{Atiyah--Singer index theorem}.~\hyperref[ch22-ref-6a]{[6a]} Among
other things, it shows that under variations in the gauge field the integral
on the left-hand side of Eq.~\eqref{eq:22.2.49} cannot change smoothly,
but only by integers, so this integral can only depend on the topology of
the gauge field. Its dependence on this topology will be described in
\hyperref[sec:23.5]{Section 23.5}.
